%Version 3.1 December 2024
%%%%%%%%%%%%%%%%%%%%%%%%%%%%%%%%%%%%%%%%%%%%%%%%%%%%%%%%%%%%%%%%%%%%%%
%%                                                                 %%
%% Please do not use \input{...} to include other tex files.       %%
%% Submit your LaTeX manuscript as one .tex document.              %%
%%                                                                 %%
%%%%%%%%%%%%%%%%%%%%%%%%%%%%%%%%%%%%%%%%%%%%%%%%%%%%%%%%%%%%%%%%%%%%%

\documentclass[pdflatex,sn-mathphys-num]{sn-jnl}

%%%% Standard Packages
\usepackage{graphicx}%
\usepackage{multirow}%
\usepackage{amsmath,amssymb,amsfonts}%
\usepackage{amsthm}%
\usepackage{mathrsfs}%
\usepackage[title]{appendix}%
\usepackage{xcolor}%
\usepackage{textcomp}%
\usepackage{manyfoot}%
\usepackage{booktabs}%
\usepackage{algorithm}%
\usepackage{algorithmicx}%
\usepackage{algpseudocode}%
\usepackage{listings}%

\raggedbottom

\begin{document}

\title[Structure-Aware Keyphrase Extraction]{Structure-Aware Keyphrase Extraction with Positional, Section, and Semantic Relevance Weighting}

\author[]{\fnm{Donggeun} \sur{Yoo}}

\author[]{\fnm{Lu} \sur{Zhang}}

\abstract{%
Keyphrase extraction from scientific documents is a fundamental task in natural language processing.
While TF-IDF provides a strong unsupervised baseline, it ignores the structural information inherent
in academic papers. We propose SAKE (Structure-Aware Keyphrase Extraction), a multiplicative scoring
framework that augments TF-IDF with three additional signals: positional weighting (phrases appearing
earlier score higher), section-aware weighting (title and abstract phrases are boosted), and semantic
relevance via Latent Dirichlet Allocation (LDA). Through an ablation study on the Krapivin benchmark
(2,304 documents), we find that section-aware weighting is the most impactful component, improving
F1@5 by 28.7\% relative to TF-IDF with positional features alone. Positional weighting contributes a
further 12.0\% improvement. However, adding LDA-based semantic relevance yields no additional gain,
suggesting that lightweight topic models are redundant in the presence of explicit structural signals.
Our best configuration (TF-IDF + Position + Section) achieves an F1@5 of 0.1539, and aggressive
section weights (title 8$\times$, abstract 4$\times$) push this to 0.1668, a 56.2\% relative
improvement over the TF-IDF baseline.}

\keywords{Keyphrase Extraction, TF-IDF, Document Structure, Positional Weighting, LDA, Ablation Study}

\maketitle

%% ================================================================
%%  1  INTRODUCTION
%% ================================================================
\section{Introduction}\label{sec:intro}

Keyphrase extraction is a core task in natural language processing that supports downstream
applications such as information retrieval, summarization, topic indexing, and scientific document
understanding. TF-IDF serves as a strong, interpretable baseline for this task because it is
efficient and easy to implement. However, TF-IDF treats a document as a bag of words and fails to
model structural information that is critical in long, structured documents like research
papers~\cite{florescu2017}.

In scientific articles, candidate phrases appearing in the title, abstract, or early sections are
typically more informative than those found only in later body text. Prior work has demonstrated that
positional information and semantic relevance can improve unsupervised keyphrase extraction
quality~\cite{florescu2017,patel2021,liang2021,ding2021,song2023}.

This project investigates whether incorporating positional and structural signals into a TF-IDF-based
framework can improve keyphrase extraction performance while maintaining interpretability, and whether
a lightweight notion of topic relevance via Latent Dirichlet Allocation
(LDA)~\cite{blei2003} can further enhance keyword ranking.

%% ================================================================
%%  2  EXPERIMENTAL PLAN
%% ================================================================
\section{Experimental Plan}\label{sec:plan}

\subsection{Working Hypotheses}\label{sec:hypotheses}

We formulate three hypotheses that our experiments aim to test:

\noindent\textbf{H1 (Structural Salience).}
Candidate phrases appearing in salient sections (title, abstract) are more informative than those in
regular body text. Assigning higher structural weights to title and abstract phrases should improve
precision in identifying primary keyphrases.

\noindent\textbf{H2 (Positional Salience).}
Phrases located earlier in the document carry higher topical weight. Applying positional weighting
based on first-occurrence position should improve recall and F1.

\noindent\textbf{H3 (Semantic Relevance).}
A lightweight topic model (LDA) can capture topical fit beyond raw frequency, further improving
ranking quality when combined with positional and section-aware features.

\subsection{Dataset}\label{sec:dataset}

We use the Krapivin Benchmark~\cite{krapivin2009}, a collection of 2,304 full-text computer
science research papers with author-assigned extractive keyphrases as ground truth. Of these, 2,283
documents have non-empty ground truth and are used for evaluation. No training split is needed since
our approach is entirely unsupervised. The dataset is loaded from HuggingFace
(\texttt{midas/krapivin}, \texttt{raw} config).

\paragraph{Preprocessing.}
The preprocessing pipeline (implemented in \texttt{ATP.ipynb}) consists of:
\begin{enumerate}
  \item \textbf{Document Parsing}: Tokens are segmented into title, abstract, and body sections
        using the \texttt{-{}-} delimiter and section identifiers (T, A, B).
  \item \textbf{Text Normalization}: Regex-based cleaning removes \LaTeX{} equations, citations,
        emails, URLs, section numbers, and special characters. Text is lowercased.
  \item \textbf{Candidate Phrase Extraction}: spaCy (\texttt{en\_core\_web\_sm}) extracts noun
        phrases using POS patterns: $(\text{Adj})^{*}(\text{Noun})^{+}$. Stopwords,
        punctuation-only strings, and excessively short/long phrases are removed.
  \item \textbf{Output}: A preprocessed pickle file (\texttt{preprocessed\_krapivin.pkl}) containing
        per-document: \texttt{doc\_id}, \texttt{title\_phrases}, \texttt{abstract\_phrases},
        \texttt{body\_phrases}, \texttt{all\_phrases} (concatenation preserving order), and
        \texttt{ground\_truth}.
\end{enumerate}

\subsection{Experimental Design}\label{sec:design}

Table~\ref{tab:fixed} lists the fixed parameters and Table~\ref{tab:free} lists the free parameters
explored in the sensitivity analysis.

\begin{table}[h]
\caption{Fixed parameters}\label{tab:fixed}
\begin{tabular*}{\textwidth}{@{\extracolsep\fill}lll@{}}
\toprule
Parameter & Value & Rationale \\
\midrule
POS pattern            & $(\text{Adj})^{*}(\text{Noun})^{+}$ & Standard practice \\
spaCy model            & \texttt{en\_core\_web\_sm}           & Lightweight, sufficient for POS \\
IDF smoothing          & $\log\!\bigl(\tfrac{N}{1+\text{df}}\bigr)$ & Standard TF-IDF \\
Evaluation cutoffs $K$ & 5, 10, 15                            & Typical keyphrase list lengths \\
Stemmer                & Porter Stemmer                       & Standard for evaluation \\
LDA topics             & 20                                   & Granularity vs.\ generalization \\
LDA passes             & 10                                   & Sufficient for convergence \\
LDA random seed        & 42                                   & Reproducibility \\
Position weight        & $\frac{1}{1+\text{pos}/\text{total}}$ & Smooth decay $[0.5,1.0]$ \\
Topic weight shift     & $0.5+\cos(\mathbf{d},\mathbf{p})$    & Modulating range $[0.5,1.5]$ \\
\botrule
\end{tabular*}
\end{table}

\begin{table}[h]
\caption{Free parameters (sensitivity analysis)}\label{tab:free}
\begin{tabular*}{\textwidth}{@{\extracolsep\fill}lll@{}}
\toprule
Parameter & Tested Values & Purpose \\
\midrule
Title weight $\alpha_T$    & 2.0, 3.0, 5.0, 8.0 & Sensitivity analysis \\
Abstract weight $\alpha_A$ & 1.5, 2.0, 3.0, 4.0 & Sensitivity analysis \\
Body weight $\alpha_B$     & 1.0 (fixed)         & Reference level \\
\botrule
\end{tabular*}
\end{table}

\subsection{Experimental Setup}\label{sec:setup}

\paragraph{Combined Scoring Model.}
For each candidate phrase $t$ in document $d$, the final score is computed as:
\begin{equation}\label{eq:score}
  \text{Score}(t) = \text{TF-IDF}(t) \;\cdot\; W_{\text{position}}(t) \;\cdot\;
                    W_{\text{structure}}(t) \;\cdot\; W_{\text{topic}}(t),
\end{equation}
where
\begin{align}
  W_{\text{position}}(t) &= \frac{1}{1 + \frac{\text{first\_pos}(t)}{\text{total\_phrases}}}, \label{eq:pos}\\[4pt]
  W_{\text{structure}}(t) &= \max\!\bigl(\alpha_B,\; [\,t \in \text{Abstract}\,]\cdot\alpha_A,\;
                             [\,t \in \text{Title}\,]\cdot\alpha_T\bigr), \label{eq:sec}\\[4pt]
  W_{\text{topic}}(t) &= 0.5 + \cos(\mathbf{d}_{\text{LDA}},\;\mathbf{t}_{\text{LDA}}). \label{eq:topic}
\end{align}

\paragraph{Ablation Study Design.}
We isolate the contribution of each feature through incremental addition:
\begin{enumerate}
  \item \textbf{TF-IDF only} --- bag-of-words frequency baseline.
  \item \textbf{TF-IDF + Position} --- adds Eq.~\eqref{eq:pos}.
  \item \textbf{TF-IDF + Position + Section} --- adds Eq.~\eqref{eq:sec}.
  \item \textbf{TF-IDF + Position + Section + Semantic (Full SAKE)} --- adds Eq.~\eqref{eq:topic}.
\end{enumerate}



\paragraph{Evaluation Metrics.}
We use Precision@$K$, Recall@$K$, and F1@$K$, with predictions matched to gold keyphrases after Porter
stemming (standard practice~\cite{florescu2017}). Results are macro-averaged across all 2,283
documents with ground truth.



%% ================================================================
%%  3  RESULTS
%% ================================================================
\section{Results}\label{sec:results}

\subsection{Ablation Study}\label{sec:ablation}

Table~\ref{tab:ablation} presents the main ablation results.

\begin{table}[h]
\caption{Ablation study --- Precision, Recall, and F1 at $K=5,10,15$. Best values per metric
in bold.}\label{tab:ablation}
\begin{tabular*}{\textwidth}{@{\extracolsep\fill}lccccccccc@{}}
\toprule
 & \multicolumn{3}{c}{$K=5$} & \multicolumn{3}{c}{$K=10$} & \multicolumn{3}{c}{$K=15$} \\
\cmidrule{2-4}\cmidrule{5-7}\cmidrule{8-10}
Model & P & R & F1 & P & R & F1 & P & R & F1 \\
\midrule
TF-IDF              & .0979 & .1334 & .1068 & .0789 & .2145 & .1104 & .0664 & .2700 & .1027 \\
+ Position          & .1093 & .1496 & .1196 & .0859 & .2340 & .1201 & .0722 & .2957 & .1117 \\
+ Pos + Sec         & \textbf{.1400} & \textbf{.1938} & \textbf{.1539} & \textbf{.1050} & \textbf{.2899} & \textbf{.1470} & \textbf{.0845} & \textbf{.3473} & \textbf{.1307} \\
Full SAKE           & .1391 & .1938 & .1534 & .1031 & .2857 & .1446 & .0835 & .3461 & .1293 \\
\botrule
\end{tabular*}
\end{table}

\paragraph{Positional weighting (H2: Supported).}
Adding position weights improves F1@5 from 0.1068 to 0.1196 (+12.0\% relative), confirming that
phrases appearing earlier in academic papers are more likely to be keyphrases.

\paragraph{Section-aware weighting (H1: Strongly Supported).}
This is the most impactful component. F1@5 jumps from 0.1196 to 0.1539 (+28.7\% relative) when
section weights are added. Title and abstract phrases receive $3\times$ and $2\times$ boosts
respectively, which dramatically improves both precision and recall. This aligns with the hypothesis
that structural salience in scientific documents provides the strongest signal for keyphrase
identification.

\paragraph{Semantic relevance via LDA (H3: Not Supported).}
Adding the LDA topic weight slightly \emph{decreases} performance (F1@5: 0.1539 $\to$ 0.1534;
F1@10: 0.1470 $\to$ 0.1446). The semantic signal from a 20-topic LDA model appears largely redundant
with position and section features, or the topic granularity may be insufficient to distinguish
keyphrases from related but non-key phrases within the same topic cluster.

\subsection{Hyperparameter Sensitivity}\label{sec:sensitivity}

Table~\ref{tab:sensitivity} shows F1 for different section weight configurations using the
TF-IDF + Position + Section model.

\begin{table}[h]
\caption{F1 at different section weight configurations}\label{tab:sensitivity}
\begin{tabular*}{\textwidth}{@{\extracolsep\fill}lccc@{}}
\toprule
Configuration & F1@5 & F1@10 & F1@15 \\
\midrule
Low: $\alpha_T{=}2.0,\;\alpha_A{=}1.5$       & .1446 & .1367 & .1252 \\
Mid: $\alpha_T{=}3.0,\;\alpha_A{=}2.0$        & .1539 & .1470 & .1307 \\
High: $\alpha_T{=}5.0,\;\alpha_A{=}3.0$       & .1625 & .1535 & .1354 \\
Very High: $\alpha_T{=}8.0,\;\alpha_A{=}4.0$  & \textbf{.1668} & \textbf{.1580} & \textbf{.1368} \\
\botrule
\end{tabular*}
\end{table}

Performance increases monotonically as section weights grow larger, suggesting the model benefits
from aggressively boosting title and abstract phrases. The ``Very High'' configuration achieves
F1@5 of 0.1668, a 56.2\% relative improvement over the TF-IDF baseline. This strong sensitivity
to section weights underscores that document structure is the dominant feature for keyphrase
extraction in scientific papers.

\subsection{Error Analysis}\label{sec:errors}

Table~\ref{tab:errors} breaks down missed keyphrases by n-gram type for the Full SAKE model at
$K{=}15$.

\begin{table}[h]
\caption{Missed keyphrases by n-gram type (Full SAKE, $K{=}15$)}\label{tab:errors}
\begin{tabular*}{\textwidth}{@{\extracolsep\fill}lcc@{}}
\toprule
N-gram Type & Missed Count & Percentage \\
\midrule
Bigram  & 3{,}534 & 61.8\% \\
Unigram & 1{,}152 & 20.2\% \\
3+ gram & 1{,}032 & 18.1\% \\
\botrule
\end{tabular*}
\end{table}

The overall hit rate at $K{=}15$ is 2{,}863/8{,}581 = 33.4\%. The majority of missed keyphrases are
bigrams (61.8\%), which is expected since bigrams form the largest category of author-assigned
keyphrases. Unigrams are missed less frequently (20.2\%), likely because single-word terms tend to
have higher TF-IDF scores. Longer n-grams (3+) account for 18.1\% of misses, often due to
exact-match failures where the candidate extraction pipeline produces slight variations of the gold
phrase.

\subsection{Qualitative Examples}\label{sec:qualitative}

Sample predictions from five randomly selected documents show that both the baseline and full model
frequently identify domain-relevant terms that are semantically close to gold keyphrases but do not
match after stemming. For instance, for document 1024923 with gold keyphrases including ``forward
simulations'' and ``backward simulations'', both models correctly retrieve ``forward simulation''
and ``backward simulation'' (matching after stemming). However, keyphrases like ``refinement
mappings'' and ``prophecy variables'' --- which are conceptual rather than frequent --- are
systematically missed.

%% ================================================================
%%  4  DISCUSSION
%% ================================================================
\section{Discussion}\label{sec:discussion}

\subsection{Summary of Findings}

Our SAKE framework demonstrates that incorporating document structure into TF-IDF-based keyphrase
extraction yields substantial improvements:

\begin{enumerate}
  \item \textbf{Section-aware weighting is the single most impactful feature}, providing
        ${\sim}29\%$ relative F1 improvement. This validates the core hypothesis that document
        structure (title, abstract, body) provides strong prior information about keyphrase salience
        in scientific papers.
  \item \textbf{Positional weighting provides consistent but modest gains} (${\sim}12\%$ relative F1
        improvement), confirming that first-occurrence position correlates with keyphrase importance.
  \item \textbf{LDA-based semantic relevance does not improve over the structure-aware model.} The
        20-topic LDA model's coarse topic assignments appear redundant with the information already
        captured by position and section features. This is a negative result, but an informative one
        --- it suggests that in the presence of explicit structural signals, lightweight topic
        modeling adds little value.
\end{enumerate}

\subsection{Challenges and Solutions}

\begin{itemize}
  \item \textbf{Noisy candidate extraction}: The spaCy POS tagger occasionally produces parsing
        artifacts. We addressed this with regex-based preprocessing to remove \LaTeX{}, citations,
        and special tokens before spaCy processing.
  \item \textbf{Hyphenated compounds}: Terms like ``graph-based'' were being split by spaCy. We
        solved this by temporarily replacing hyphens with underscores during POS tagging, then
        restoring them.
  \item \textbf{Exact-match evaluation}: Stemming-based matching misses near-synonyms (e.g.,
        ``automaton'' vs.\ ``automata''). This is a known limitation of extractive keyphrase
        evaluation.
\end{itemize}

\subsection{Limitations}

\begin{itemize}
  \item The approach is limited to extractive keyphrases --- it cannot generate abstractive
        keyphrases that do not appear verbatim in the text.
  \item The LDA model uses a fixed number of topics (20), which may not be optimal. Grid search
        over topic counts could potentially improve the semantic component.
  \item Section boundary detection relies on dataset-specific markers (\texttt{-{}-T}, \texttt{-{}-A},
        \texttt{-{}-B}). Generalizing to arbitrary documents would require a robust section detector.
\end{itemize}

\subsection{Future Directions}

\begin{enumerate}
  \item \textbf{Better semantic representations}: Replace LDA with sentence-level embeddings (e.g.,
        Sentence-BERT) for computing phrase--document similarity.
  \item \textbf{Learned weight combination}: Instead of fixed multiplicative weights, use a small
        held-out set to learn optimal combination weights via logistic regression or grid search.
  \item \textbf{Cross-dataset evaluation}: Test on additional keyphrase benchmarks (SemEval-2010,
        Inspec, NUS) to assess generalizability.
  \item \textbf{Abstractive keyphrase generation}: Explore hybrid models that combine extractive
        candidate ranking with generative keyphrase prediction.
\end{enumerate}

%% ================================================================
%%  5  CONCLUSION
%% ================================================================
\section{Conclusion}\label{sec:conclusion}

We presented SAKE, a simple yet effective unsupervised framework for keyphrase extraction from
scientific documents that augments TF-IDF with positional, section-aware, and semantic relevance
signals through multiplicative scoring. Our ablation study on the Krapivin benchmark (2,304
documents) reveals a clear hierarchy of feature importance: section-aware weighting contributes the
largest gain (+28.7\% relative F1@5), followed by positional weighting (+12.0\%), while LDA-based
semantic relevance provides no additional benefit. The best-performing configuration (TF-IDF +
Position + Section with aggressive title/abstract weights) achieves an F1@5 of 0.1668, representing
a 56.2\% relative improvement over the TF-IDF baseline.

These findings highlight that explicit structural signals in scientific papers --- particularly the
distinction between title, abstract, and body --- are the dominant cue for keyphrase salience in
unsupervised settings. Our negative result on LDA suggests that coarse-grained topic models are
redundant when structural features are already present, pointing toward richer semantic
representations (e.g., contextual embeddings) as the more promising direction for future improvement.


%% ================================================================
%%  6  TEAM RESPONSIBILITIES
%% ================================================================
\section{Team Responsibilities}\label{sec:team}

\begin{itemize}
  \item \textbf{Donggeun Yoo}: Idea proposing, dataset parsing, text preprocessing, candidate
        phrase extraction, and implementation of the TF-IDF baseline.
  \item \textbf{Lu Zhang}: Implementation of the positional, section-aware, and semantic relevance
        modules; conducting ablation experiments; writing the report.
\end{itemize}

%% ================================================================
%%  BACKMATTER
%% ================================================================
\backmatter

\bmhead{Acknowledgements}

This work was completed as part of the CSC 8855 course project.

\section*{Declarations}

\begin{itemize}
\item \textbf{Code availability}: All code is available in the companion Jupyter notebooks
      \texttt{ATP.ipynb} and \texttt{SAKE\_Methodology.ipynb}.
\item \textbf{Data availability}: The Krapivin dataset is publicly available via HuggingFace
      (\texttt{midas/krapivin}).
\end{itemize}

\bibliography{sn-bibliography}

\end{document}
